\section{Counterexample search and adversarial diagnostics}
\label{cx:section}

\noindent\textbf{Scope: diagnostics and model obstructions.}
The actual fixed-section counterexample appears later as
Theorem~\ref{uc:counterexample}.
The constructions below separate exact Maxwell counterexamples to proposed
\emph{intermediate} lemmas from exact finite-dimensional models. No model port
is identified with a physical channel through the prescribed section without
an additional proof. The material class throughout the physical discussion is
the fixed, real, scalar class $1\leq\varepsilon_r\leq\varepsilon_+$ at one
positive frequency, with vacuum exterior. The section, frequency, and
efficiency and bypass thresholds do not vary. Support size and resonance
lifetime may grow. The radial, local-flux, and mirror-network diagnostics
originate in the prior report \cite{priorreport}; they are rederived here
because their precise scope is consequential.

\subsection{An exact bounded-index Maxwell amplification obstruction}

\begin{proposition}[PROVED: unbounded local amplification, one reflected channel]
\label{cx:bragg}
For every fixed $\varepsilon_+>1$ and $k_0>0$ there are a fixed ball
$B_{r_0}$ and finite radial dielectric structures $\varepsilon_M$, with
$1\leq\varepsilon_M\leq\varepsilon_+$, for which a unit-power incoming
spherical radiation channel produces electric and magnetic squared field
norms in $B_{r_0}$ growing proportionally to $q^{2M}$, for some fixed $q>1$.
The exterior radius grows linearly with $M$. The selected scattering
coefficient has modulus one and the net power through every concentric
sphere is zero.
\end{proposition}

\begin{proof}
Choose $1<q<\sqrt{\varepsilon_+}$. Write
$\widehat j_1(x)=xj_1(x)=\sin(x)/x-\cos(x)$. Its derivative has arbitrarily
large positive roots: at $(m+1/2)\pi>1$ its sign alternates, since its value
is $(-1)^m(1-x^{-2})$. Choose such a root $x_0>1$ large enough that
$q^2+2/x_0^2\leq\varepsilon_+$, and put $r_0=x_0/k_0$.
The root is not simultaneously a zero of $\widehat j_1$, by uniqueness for
its second-order ordinary differential equation away from zero.

For a tangential, surface-divergence-free vector spherical harmonic
$X_{1m}$ with $\int_{S^2}|X_{1m}|^2\,d\Omega=1$, set
\begin{equation}
 E=\frac{u(r)}r X_{1m},\qquad
 H=\frac1{i\omega\mu_0}\nabla\times E,
 \qquad
 u''+\left(k_0^2\varepsilon_M(r)-\frac2{r^2}\right)u=0.
 \label{cx:radial-maxwell}
\end{equation}
This is the exact TE multipole reduction of both Maxwell curl equations.
At radial interfaces continuity of tangential $E$ and $H$ is exactly
continuity of $u,u'$; permeability is constant and the normal electric
component vanishes. The integrated outgoing power is
\begin{equation}
 P(r)=\frac1{2\omega\mu_0}\operatorname{Im}(\overline u u').
 \label{cx:radial-power}
\end{equation}
In the core take $\varepsilon_M=1$ and
$u=C_M\widehat j_1(k_0r)$, which is regular at zero and has $u'(r_0)=0$.

Append $M$ low/high shell pairs. A shell of type $n\in\{1,q\}$ has
width $\pi/(2k_0n)$ and coefficient
$\varepsilon_M(r)=n^2+2/(k_0r)^2$. Thus its radial equation is exactly
$u''+(k_0n)^2u=0$. On the state $(u,u'/k_0)^T$, its transfer matrix is
$\left(\begin{smallmatrix}0&1/n\\-n&0\end{smallmatrix}\right)$.
A pair has transfer matrix $-\operatorname{diag}(q^{-1},q)$, so
\begin{equation}
 R_M=r_0+\frac{M\pi}{2k_0}(1+q^{-1}),\quad
 u(R_M)=(-q^{-1})^M C_M\widehat j_1(x_0),\quad u'(R_M)=0.
 \label{cx:bragg-transfer}
\end{equation}
All shells obey the claimed material bounds. In the vacuum exterior put
$u=a\widehat h^{(2)}_1(k_0r)+b_M\widehat h^{(1)}_1(k_0r)$, where
$\widehat h^{(1)}_1(x)=-e^{ix}(1+i/x)$ and
$\widehat h^{(2)}_1=\overline{\widehat h^{(1)}_1}$ on the positive real
axis. Direct differentiation gives the Wronskian $2i$ and
\begin{equation}
 |(\widehat h^{(1)}_1)'(x)|^2=1-x^{-2}+x^{-4},\quad
 b_M=-a\frac{(\widehat h^{(2)}_1)'(k_0R_M)}
                    {(\widehat h^{(1)}_1)'(k_0R_M)},\quad
 u(R_M)=\frac{2ia}{(\widehat h^{(1)}_1)'(k_0R_M)}.
 \label{cx:bragg-exterior}
\end{equation}
Incoming power is $k_0|a|^2/(2\omega\mu_0)$; choose
$|a|^2=2\omega\mu_0/k_0$. Consequently
\begin{equation}
 |C_M|^2=
 \frac{4q^{2M}|a|^2}{|\widehat j_1(x_0)|^2
             |(\widehat h^{(1)}_1)'(k_0R_M)|^2}.
 \label{cx:bragg-growth}
\end{equation}
For $x\geq1$ the derivative factor squared lies in $[3/4,1]$.
The core fields are fixed nonzero profiles times $C_M$, proving the
norm statement. A fixed interior section on which their traces are
nonzero has the corresponding trace growth. Equation~\eqref{cx:bragg-exterior}
gives $|b_M|=|a|$, while the core radial function is a constant times a
real function. Equation~\eqref{cx:radial-power} and conservation then give
zero total radial flux at every radius. No independently chosen trapped
amplitude was inserted: $C_M$ is fixed by the incoming amplitude and
the interface matching.
\end{proof}

\begin{remark}[DISPROVED shortcut; scope of the example]
\label{cx:bragg-scope}
External passivity does not uniformly bound the incident-power-to-local-field
map over the admitted class with unrestricted support. The prior report's
$q=2$ example required $\varepsilon_+>4$; choosing $x_0$ and $q$ as above
extends this obstruction to every $\varepsilon_+>1$. This does not disprove
a uniform channel count, a fixed-support estimate, or a theorem restricted
to counted fields with positive through-power. A complete spherical
scattering phase is not an additional transmitted channel. This exact
construction establishes neither an increasing number of transmitted
directions nor the required absolute bypass bound.
\end{remark}

\subsection{Evanescent pairs can carry real power locally}

\begin{proposition}[PROVED: infinite positive local flux index]
\label{cx:evanescent}
On a fixed rectangular section $D=(0,L)\times(0,W)$ and homogeneous
source-free collar $D\times(-h,h)$, the space of smooth local Maxwell
solutions has positive flux subspaces of arbitrarily large finite
dimension. This remains true if every superposition has zero real
normal flux on the four collar sidewalls.
\end{proposition}

\begin{proof}
Let $k^2=\omega^2\mu_0\varepsilon_0\varepsilon_c$ and take every integer
$m$ with $q_m=m\pi/L>k$. Put $\kappa_m=\sqrt{q_m^2-k^2}$ and
\begin{equation}
 f_m(z)=a_m(e^{-\kappa_mz}+i e^{\kappa_mz}),\qquad
 |a_m|^2=\frac{2\omega\mu_0}{\kappa_mLW},\qquad
 E_m=(0,f_m\sin(q_mx),0),
 \label{cx:evanescent-electric}
\end{equation}
\begin{equation}
 H_m=\frac1{\omega\mu_0}
       (i f'_m\sin(q_mx),0,-i q_m f_m\cos(q_mx)).
 \label{cx:evanescent-magnetic}
\end{equation}
Both curl equations hold by $f_m''=(q_m^2-k^2)f_m$.
The normal power density is
\begin{equation}
\frac{\operatorname{Im}(\overline f_m f'_m)\sin^2(q_mx)}{2\omega\mu_0}
=\frac{\kappa_m|a_m|^2\sin^2(q_mx)}{\omega\mu_0}.
\label{cx:evanescent-density}
\end{equation}
Sine orthogonality therefore gives, for every finite coefficient vector,
\begin{equation}
 P_D\left(\sum_m c_m(E_m,H_m)\right)=\sum_m|c_m|^2.
 \label{cx:evanescent-gram}
\end{equation}
On the $x$ sidewalls $E$ vanishes, and on the $y$ sidewalls the normal
Poynting component vanishes identically. Both statements hold for
every finite superposition, not just for basis fields.
\end{proof}

\begin{remark}[DISPROVED local cutoff; OPEN global completion]
\label{cx:evanescent-scope}
A single decreasing exponential carries zero real normal power. Its
coherent combination with an increasing exponential does not. Thus a
propagating-wave count cannot be justified by discarding each evanescent
term separately. These are exact \emph{local} vector Maxwell fields in
ordinary material, but their norms at the collar ends grow rapidly with
transverse order. No finite passive dielectric device producing these
data from power-normalized external channels has been constructed.
Zero sidewall flux is a property of the fields, not permission to install
conducting walls. Nor does it control the field on the rest of the
infinite section plane. Hence this proposition does not establish
global absolute bypass or a physical counterexample.
\end{remark}

\subsection{An exact feedback diagnostic independent of mode reduction}

\begin{proposition}[PROVED in lossless scattering algebra]
\label{cx:mirror-network}
For any $N$ and any $0<\xi_j\leq1$, reciprocal unitary two-port networks
with isolated transmission powers $\xi_j$ can be cascaded with lossless
phases to produce $N$ perfectly transmitted orthogonal network channels.
\end{proposition}

\begin{proof}
For direction $j$, use two identical mirrors with
$r_j=i\sqrt{1-\xi_j}$ and $t_j=\sqrt{\xi_j}$. The symmetric matrix
$\left(\begin{smallmatrix}r_j&t_j\\t_j&r_j\end{smallmatrix}\right)$
is unitary. With propagation phase $p_j=i$, the convergent reflection
series gives
\begin{equation}
 t_{{\rm all},j}=\frac{t_j^2p_j}{1-r_j^2p_j^2}=i,
 \qquad
 r_{{\rm all},j}=r_j+
       \frac{t_j^2r_jp_j^2}{1-r_j^2p_j^2}=0.
 \label{cx:mirror-feedback}
\end{equation}
The direct sum has transmission $iI_N$. The feedback inverse is exactly
$\operatorname{diag}(\xi_j^{-1})$, compensating arbitrarily weak isolated
links.
\end{proof}

\noindent\textbf{What is disproved.}
An isolated-link singular-value tail together with external passivity
does not control the completed scattering transmission. This proposition
is stronger as an algebraic diagnostic than a coupled-mode approximation,
but still supplies no bounded-index synthesis through the same physical $D$.

\subsection{Matched resonator pairs: energy, bandwidth, and errors}

The following model uses the energy/power normalization conventional in
temporal coupled-mode theory \cite{fan2003,suh2004}. Its equations and all
conclusions about them are proved directly; their agreement with a
particular Maxwell device remains an additional question.

\begin{definition}[MODEL: a reciprocal resonator pair]
\label{cx:pair-model}
Let $x,y$ be energy-normalized resonator amplitudes, so stored model energy
is $|x|^2+|y|^2$. Let $a,b$ be incoming amplitudes with incident power
$|a|^2+|b|^2$. For $J,\gamma>0$, set
\begin{align}
 \dot x&=(-i\omega_0-\gamma)x-iJy+\sqrt{2\gamma}\,a,
 &a_-&=-a+\sqrt{2\gamma}\,x,\notag\\
 \dot y&=(-i\omega_0-\gamma)y-iJx+\sqrt{2\gamma}\,b,
 &b_-&=-b+\sqrt{2\gamma}\,y.
 \label{cx:pair-dynamics}
\end{align}
The $-1$ direct reflection is part of the model's port convention.
\end{definition}

\begin{proposition}[PROVED in the model: exact matching and cost]
\label{cx:pair-solution}
For harmonic drive at $\omega=\omega_0+\Delta$ with $a=1,b=0$, define
$D_\Delta=(\gamma-i\Delta)^2+J^2$. Then
\begin{align}
 x&=\frac{\sqrt{2\gamma}(\gamma-i\Delta)}{D_\Delta},&
 y&=\frac{-i\sqrt{2\gamma}J}{D_\Delta},\notag\\
 t&=\frac{-2i\gamma J}{D_\Delta},&
 r&=\frac{\gamma^2+\Delta^2-J^2}{D_\Delta}.
 \label{cx:pair-scattering}
\end{align}
The model is lossless and passive. At $\Delta=0$, $\gamma=J$ it has
$t=-i,r=0$, internal energy $J^{-1}$, and left-to-right internal current
$1$. At this matched setting,
\begin{equation}
 |t(\Delta)|^2=\frac1{1+\Delta^4/(4J^4)},\qquad
 {\rm FWHM}=2\sqrt2 J.
 \label{cx:matched-linewidth}
\end{equation}
\end{proposition}

\begin{proof}
Substitution of harmonic time dependence gives the linear system with
matrix $\left(\begin{smallmatrix}\gamma-i\Delta&iJ\\
iJ&\gamma-i\Delta\end{smallmatrix}\right)$; inversion proves
\eqref{cx:pair-scattering}. Direct differentiation of the energy gives
\begin{equation}
 \frac{d}{dt}(|x|^2+|y|^2)
 =|a|^2+|b|^2-|a_-|^2-|b_-|^2.
 \label{cx:pair-energy-balance}
\end{equation}
The internal left-to-right exchange is
$I_{L\to R}=-2J\operatorname{Im}(\overline x y)$, from the individual
energy derivatives. At matching $x=1/\sqrt{2J}$ and $y=-ix$, giving the
energy and current. Finally,
$|D_\Delta|^2=(2J^2-\Delta^2)^2+4J^2\Delta^2
=4J^4+\Delta^4$ gives the line shape and its half-power frequencies.
The poles are at $\Delta=\pm J-i\gamma$, and hence remain strictly
below the real axis for every member of the family.
\end{proof}

\begin{remark}[PROVED model sensitivities]
\label{cx:pair-sensitivity}
At $\Delta=0$, writing $s=\gamma/J$ gives
$|t|^2=4s^2/(1+s^2)^2$. Thus $|t|^2\geq\tau$ exactly when
\begin{equation}
 \frac{1-\sqrt{1-\tau}}{\sqrt\tau}\leq s\leq
 \frac{1+\sqrt{1-\tau}}{\sqrt\tau}.
 \label{cx:rate-tolerance}
\end{equation}
At matching, a common frequency detuning obeys the threshold exactly when
$|\Delta|\leq\sqrt2 J(\tau^{-1}-1)^{1/4}$.
If both resonators instead have an additional unobserved amplitude decay
$\kappa\geq0$, replace the diagonal rate in the equations by
$\gamma+\kappa$ while keeping the driven/output couplings
$\sqrt{2\gamma}$. For $\gamma=J$ and $\Delta=0$,
\begin{equation}
 |t|^2=\frac4{[(1+\kappa/J)^2+1]^2}.
 \label{cx:loss-tolerance}
\end{equation}
Consequently fixed loss destroys the small-$J$ limit. To approach unit
transmission it is necessary here that $\kappa/J\to0$.
For a fixed threshold, the exact permitted ratio is
$\kappa/J\leq\sqrt{2/\sqrt\tau-1}-1$.
These costs cannot be suppressed by saying merely that each cavity has
large $Q$: its leakage and errors must be compared to its weak link.
\end{remark}

\begin{theorem}[PROVED in the model: $N$ weak directions and a stable certificate]
\label{cx:n-model}
Let $J=\operatorname{diag}(J_1,\ldots,J_N)$ with all $J_j>0$. At the
common frequency $\omega_0$, take $N$ independent pairs with external
rates $J_j$. Put
\begin{equation}
 G=\operatorname{diag}(J^{1/2},J^{1/2}),\quad
 B_0=\begin{pmatrix}I&iI\\iI&I\end{pmatrix},\quad
 A_0=GB_0G,\quad
 E_L a=(a,0)^T,\quad E_R a=(0,a)^T.
 \label{cx:weighted-matrices}
\end{equation}
Their transmission and reflection are $T_0=-iI_N$, $R_0=0$; their
internal energy matrix is $J^{-1}$ and exchange-current matrix is $I_N$.
If the system matrix changes to $A=A_0+\mathcal E$, keeping the same
port couplings, define
$F=G^{-1}\mathcal E G^{-1}$ and $\rho=\|F\|$. Whenever
$\rho<\sqrt2$, the system remains invertible and
\begin{equation}
 \|T-T_0\|\leq\frac{\rho}{1-\rho/\sqrt2},\qquad
 T=2E_R^*(B_0+F)^{-1}E_L.
 \label{cx:weighted-error}
\end{equation}
In particular every input is transmitted with power greater than
$\tau$ if
\begin{equation}
 \rho<\frac{1-\sqrt\tau}{1+(1-\sqrt\tau)/\sqrt2}.
 \label{cx:rho-certificate}
\end{equation}
\end{theorem}

\begin{proof}
The unperturbed statement follows pairwise from
Proposition~\ref{cx:pair-solution}. Also $B_0^*B_0=2I$, so
$\|B_0^{-1}\|=1/\sqrt2$ and
$\|(B_0+F)^{-1}\|\leq1/(\sqrt2-\rho)$ by the Neumann series.
The resolvent identity gives
\begin{equation}
 \|T-T_0\|
 \leq2\|(B_0+F)^{-1}\|\|F\|\|B_0^{-1}\|
 \leq\frac{\sqrt2\rho}{\sqrt2-\rho}.
 \label{cx:resolvent-proof}
\end{equation}
The triangle inequality gives $\|Ta\|\geq(1-\|T-T_0\|)\|a\|$.
Solving the sufficient strict inequality
$\rho/(1-\rho/\sqrt2)<1-\sqrt\tau$ gives
\eqref{cx:rho-certificate}.
\end{proof}

\begin{remark}[Model generality and port errors]
\label{cx:port-errors}
The perturbation may mix different pairs and need not be diagonal.
A sufficient unweighted condition is
$\|\mathcal E\|/\min_jJ_j\to0$, since
$\rho\leq\|\mathcal E\|/\min_jJ_j$; weighted control can be much sharper.
If normalized input/output couplings also change to $E_L+L$ and $E_R+R$,
and there is a direct transmission correction $D_0$, the same proof yields
\begin{equation}
 \|T-T_0\|\leq
 \frac{\rho}{1-\rho/\sqrt2}
 +\frac{2(\|L\|+\|R\|+\|L\|\|R\|)}{\sqrt2-\rho}
 +\|D_0\|.
 \label{cx:all-errors}
\end{equation}
The normalization here is relative to the original $G$; for example the
physical input matrix in the model is $\sqrt2G(E_L+L)$.
An arbitrary algebraic perturbation is not automatically passive or
reciprocal. Any proposed physical realization must also establish those
properties and the power normalization of its actual ports. The estimate
is a sufficient matrix certificate once those facts hold.
\end{remark}

\subsection{Compression and coherent passage certification}

\begin{proposition}[PROVED: physical channel certificate if its data exist]
\label{cx:compression}
Let $T:\mathcal H_{\rm in}\to\mathcal H_{\rm out}$ be a bounded physical
transmission operator between power Hilbert spaces. Let $V$ and $Z$ be
isometries from $\mathbb C^N$ into its input and output spaces, and let
$M=Z^*TV$. If $U$ is unitary and
$\|M-U\|\leq e<1-\sqrt\tau$, then
$\|TVa\|^2>\tau\|a\|^2$ for every nonzero $a$.
If $T$ is a contraction, and its physical section-flux matrix on these
inputs obeys $Q_D\succeq\eta I_N$, then the passage inequality
$P_D(Va)\geq\eta\|TVa\|^2$ also holds for every $a$.
\end{proposition}

\begin{proof}
The isometry $Z$ gives $\|TVa\|\geq\|Z^*TVa\|$; the unitary comparison
gives $\|Ma\|\geq(1-e)\|a\|$. For the second statement,
$P_D(Va)=a^*Q_Da\geq\eta\|a\|^2\geq\eta\|TVa\|^2$.
The exact, potentially weaker passage test is
$Q_D-\eta V^*T^*TV\succeq0$. None of these conclusions needs compactness
or a discrete singular-value decomposition of the full operator.
\end{proof}

\begin{lemma}[PROVED: a sufficient absolute-bypass certificate]
\label{cx:absolute-bypass}
Suppose actual basis fields have well-defined pointwise traces on
$\Pi\setminus D$, where $\Pi$ is the section plane, and define the
Hermitian matrix
\begin{equation}
 F_{ij}(s)=\frac14\left(E_j(s)\times\overline{H_i(s)}
                 +\overline{E_i(s)}\times H_j(s)\right)\cdot n.
 \label{cx:pointwise-flux}
\end{equation}
If these traces are regular enough that $\int_{\Pi\setminus D}\|F(s)\|\,dA$
is defined and at most $\delta$, then every coherent vector satisfies
\begin{equation}
 \int_{\Pi\setminus D}|S_n(a;s)|\,dA\leq\delta\|a\|^2.
 \label{cx:bypass-certificate}
\end{equation}
\end{lemma}

\begin{proof}
Expansion gives $S_n(a;s)=a^*F(s)a$. Pointwise,
$|a^*F(s)a|\leq\|F(s)\|\|a\|^2$; integrate.
\end{proof}

\begin{remark}[DISPROVED basis-only and signed-flux tests]
\label{cx:coherent-failure}
Let $F(s)=h(s)\left(\begin{smallmatrix}0&L\\L&0\end{smallmatrix}\right)$,
with $\int|h|=1$. Each coordinate basis vector has zero local flux,
but $a=(1,1)^T/\sqrt2$ has absolute integrated flux $L$. If additionally
$\int h=0$, the full \emph{signed} integrated flux matrix vanishes while
the same absolute bypass remains $L$. This elementary example disproves
both inference rules for Hermitian flux data; it is not asserted to be
a global Maxwell realization. The evanescent pair in
Proposition~\ref{cx:evanescent} gives an actual local Maxwell instance
of individually zero-flux exponential fields acquiring power by
interference. Likewise a basis test for signed passage can fail:
$Q=\left(\begin{smallmatrix}1&2\\2&1\end{smallmatrix}\right)$ has positive
diagonal but a negative eigenvalue. Section-flux conditions require a
full matrix inequality; absolute bypass requires additional control such
as Lemma~\ref{cx:absolute-bypass}. Rough Maxwell traces are not silently
assumed to have these pointwise products.
\end{remark}

\subsection{Physical realization attempts and their specific obstructions}

\begin{proposition}[PROVED: a restricted radial angular-momentum obstruction]
\label{cx:angular-obstruction}
For a radial scalar dielectric with $\varepsilon(r)\leq\varepsilon_+$
on $0<r\leq R$, a regular real TE radial solution of angular order
$\ell\geq1$ obeys
\begin{equation}
 u''=\left(\frac{\ell(\ell+1)}{r^2}-k_0^2\varepsilon(r)\right)u.
 \label{cx:angular-equation}
\end{equation}
If $\ell(\ell+1)>k_0^2\varepsilon_+R^2$ and the nonzero leading
coefficient in $u\sim cr^{\ell+1}$ is positive, then $u,u'>0$ on
$(0,R]$. In particular, neither $u(R)=0$ nor $u'(R)=0$ can hold.
\end{proposition}

\begin{proof}
The regular solution also satisfies
$u'(r)\sim(\ell+1)c r^\ell$. To justify this from the equation rather
than differentiate an asymptotic without control, put
$v=r^{-\ell-1}u$. Then
$(r^{2\ell+2}v')'=-k_0^2\varepsilon(r)r^{2\ell+2}v$.
The integration constant at zero must vanish for $v$ to have its finite
limit $c$, so boundedness of $\varepsilon,v$ gives $v'=O(r)$.
Thus both $u$ and $u'$ are positive for sufficiently small positive radius.
As long as $u>0$, the displayed strict inequality gives $u''>0$, so
$u'$ is increasing and remains positive. A first zero of either $u$
or $u'$ is consequently impossible. Across radial interfaces $u,u'$
are continuous; the same argument applies piecewise and hence globally.
\end{proof}

\noindent\textbf{Radial shells and partial waves.}
Proposition~\ref{cx:bragg} proves that bounded-index reflectors can produce
arbitrarily large internal amplitude. Repeating its exactly compensated
shell profile for order $\ell$ would require
$\varepsilon=n^2+\ell(\ell+1)/(k_0r)^2$. At a fixed inner radius this
eventually violates the fixed material upper bound. Increasing the
radius avoids that violation but does not place the high-order power
through the original $D$. Proposition~\ref{cx:angular-obstruction}
also blocks the particular real-core Neumann or Dirichlet tuning in a
fixed small ball at high $\ell$. It is not a general exclusion of open
resonances, nonspherical matching, or fields coupled from a larger
exterior structure. The exact radial example is a one-port reflector;
turning its stored energy into $N$ independently driven opposite-side
ports is an unproved step.

\medskip
\noindent\textbf{Planar dielectric Bragg cavities.}
At normal incidence, bounded-index quarter-wave layers do realize
lossless one-dimensional mirrors; a symmetric resonant pair can realize
the scalar feedback algebra. This statement follows by the same
two-by-two transfer-matrix multiplication used above. Such an exact
planar realization has infinite transverse extent and its plane waves
are normalized per transverse area. It does not supply finite total
incoming power with absolute flux confined to a prescribed finite $D$.
Truncating the mirrors adds lateral leakage and diffraction. An $N$-cavity
proposal must prove that these errors, along with off-diagonal mixing and
detuning, satisfy the relative bounds of
Theorem~\ref{cx:n-model} for the weakest direction. Merely increasing
the number of mirror layers supplies no such multidimensional estimate.

\medskip
\noindent\textbf{Many small tips coupled to large remote resonators.}
Arbitrarily fine dielectric tips or corners could create linearly
independent local trace patterns. Linear independence does not bound
the smallest singular value of the physical coupling matrix, and the
field-to-output inverse is precisely where its weakness might be
compensated. No bounded-index synthesis is established that makes all
$N$ resonances coincide at the prescribed frequency while controlling
leakage, parasitic coupling, and model error relative to that smallest
coupling. Remote support growth is allowed, so its expense alone is not
an obstruction. In contrast, assuming ideal independently tunable
impedances or perfect conducting partitions would change the physical
class. These routes remain candidates, not impossibility results.

\medskip
\noindent\textbf{Thin dielectric guides and broadened channels.}
A weakly confined guide mode, when it exists, does not automatically
satisfy the fixed-aperture bypass requirement. The prior report's
qualified planar-guide construction explicitly changes the lower
permittivity bound toward zero, varies the section shape with unbounded
diameter, and uses infinite leads. It therefore cannot be transferred
to the present fixed ordinary-dielectric, fixed-background, same-section
problem. Counting outputs at different frequencies, adding internal
sources, or sending the threshold $\tau$ to zero would likewise answer
different questions.

\begin{remark}[OPEN: the missing physical bridge, stated quantitatively]
\label{cx:missing-bridge}
For each $N$ one would need an actual finite dielectric and power
isometries $V_N,Z_N$ at the same frequency whose completed compression
obeys the bound in \eqref{cx:all-errors} with a fixed margin below
$1-\sqrt\tau$, together with the physical section inequality in
Proposition~\ref{cx:compression} and an absolute-bypass estimate such
as \eqref{cx:bypass-certificate}. In a resonator implementation this
requires a justified Maxwell reduction and a weighted error matrix;
$o(\min_jJ_j)$ unweighted error is sufficient but has not been proved.
The matrices $J^{-1}$ and $I_N$ in Theorem~\ref{cx:n-model} are model
energy and internal-current matrices, not computed Maxwell energy or
section-flux matrices. If this bridge were established for arbitrary
$N$, the fixed-parameter finite channel bound would be false. Its
absence means that neither a disproof nor a universal bound has been
obtained by this track.
\end{remark}
